\begin{resumo}[Abstract]
\begin{otherlanguage*}{english}
\begin{SingleSpace}
This study investigates the extent to which the type and order of perspectives assigned to agents based on large language models alter the content and affective tone of news texts across rewriting chains. To measure these transformations, it proposes the \textit{Structured Topic Drift Index} (STDI), which decomposes drift into theme, subtopics, entities, and relations, with increments for internal contradiction and affective change. The study comprises six final simulation runs, combining three generative models and two corpora of 50 entries each: political news descriptions in English and captions from journalistic posts in Portuguese in the Brazilian electoral context. Each run applies eight trajectories of four steps, organized into trees with shared initial steps, yielding a total of 7,200 unique rewrites. The analysis combines quantitative comparisons paired by input, qualitative inspections, and estimates of valence, arousal, and dominance, distinguishing distance from the original, incremental change, and the sum of increments. The results show that rewrites can preserve the central event and its participants while altering justifications, motives, and relations. Subtopics and relations have the highest mean scores among the content components. Neutral relay from the outset is associated with greater preservation, whereas later neutral steps retain received interpretations without automatically recovering the original. Trajectories involving a conspiratorial perspective exhibit, on average, lower valence and higher arousal according to the language model used for affective assessment. The comparisons do not identify an ordering that consistently produces greater drift across generators and corpora. Disagreements between methods and score fluctuations in equivalent texts also reveal measurement limitations. The study concludes that analyzing trajectories and components makes it possible to locate transformations that a final score alone cannot explain. Its contributions are instrumental and descriptive, constrained by corpus selection, the common evaluator, and the absence of an independent numerical human reference. The STDI describes textual changes without establishing external factual veracity or measuring human behavior in news circulation.
\end{SingleSpace}

\vspace{\onelineskip}
\textbf{Keywords:} informational drift; artificial intelligence (AI); large language models (LLM); computational agents; news rewriting; affective analysis.

\end{otherlanguage*}
\end{resumo}
